% T5 summary (compact-by-compact and limit)
\subsection{T5 (legacy): compact-by-compact transfer skeleton}

\noindent\textbf{Status note.} This subsection is retained only as an archival record of the old compact-transfer strategy. It is not part of the current proof because it implicitly assumed that the centered positivity cone from A3 was already dense in each $W_K$; the main text now isolates that missing compatibility step explicitly in Section~\ref{sec:Main-closure}.
\begin{definition}[Weil inductive-limit topology]\label{t5:def:Weil_LF}
Let $W_K:=C^+_{\mathrm{even}}([-K,K])$ with the uniform norm. Define the Weil class $W:=\bigcup_{K\ge1} W_K$ with the inductive (LF) topology: $U\subset W$ is open iff $U\cap W_K$ is open in $W_K$ for every $K$. A quadratic functional $Q^\star(t;\cdot):W\to\RR$ is (sequentially) continuous in this topology iff each restriction $Q^\star(t;\cdot)|_{W_K}$ is continuous in $\|\cdot\|_\infty$.
\end{definition}

\begin{lemma}[Local continuity suffices for T5]\label{t5:lem:local-Lip}
If for every $K$ the restriction $Q^\star(t;\cdot)|_{W_K}$ is Lipschitz in $\|\cdot\|_\infty$ with some (possibly $K$-dependent) constant $L_K$, then the inductive-limit topology of Definition~\ref{t5:def:Weil_LF} guarantees sequential continuity of $Q^\star(t;\cdot)$ on $W$. No uniform bound $\sup_K L_K<\infty$ is required: whenever $\Phi_n\to\Phi$ in $W$, the convergence takes place in a single $W_K$, and the corresponding $L_K$ controls $|Q^\star(t;\Phi_n)-Q^\star(t;\Phi)|$.
\end{lemma}

\begin{lemma}[T5: transfer across $K\uparrow$]\label{t5:lem:T5_transfer}
If $Q^\star(t;\Phi)\ge0$ on every $W_K$ and the family $\{Q^\star(t;\cdot)|_{W_K}\}$ is compatible with the natural inclusions $W_K\hookrightarrow W_{K'}$ for $K<K'$, then $Q^\star(t;\Phi)\ge0$ on $W$.
\end{lemma}

\begin{proposition}[LF--transfer of positivity]\label{t5:prop:T5_transfer_formal_aux}
Let $\{W_K\}_{K\in\mathbb N}$ be an increasing family of cones of even, nonnegative $C_c$ tests supported in $[-K,K]$, and let $W=\varinjlim W_K$ be their LF inductive limit. Suppose: (i) for each $K$, the quadratic form $Q^\star(t;\cdot)$ is continuous on $W_K$ in the $\|\cdot\|_\infty$ topology; (ii) $Q^\star(t;\Phi)\ge 0$ for all $\Phi\in W_K$ for every $K$; and (iii) the embeddings $W_K\hookrightarrow W_{K+1}$ are continuous and compatible with $Q^\star(t;\cdot)$. Then $Q^\star(t;\Phi)\ge 0$ on $W$.
\end{proposition}

\begin{remark}
Continuity in (i) uses the local constants $L_K$ from Corollary~\ref{cor:A2-Lip}. We never require a uniform bound in $K$: the inductive-limit topology only asks for continuity on each fixed $W_K$, which is provided by A2.
\end{remark}

\begin{remark}[Uniform parameters]\label{rem:T5-uniform}
The uniform Archimedean floor $c_*>0$ from Lemma~\ref{lem:uniform-arch-floor}
and the uniform RKHS cap $\rho(t)\le c_*/4$ from Corollary~\ref{cor:uniform-prime-cap}
provide K-independent bounds that work for all compacts simultaneously.
The inductive-limit transfer of Proposition~\ref{t5:prop:T5_transfer_formal_aux}
would propagate positivity from the uniform A3 bridge to the full Weil class once a common dense positive cone were available.
\end{remark}

\begin{proof}
Given $\Phi\in W$, pick $K$ with $\operatorname{supp}\Phi\subset[-K,K]$; then $\Phi\in W_K$ and $Q^\star(t;\Phi)\ge0$ by (ii). Compatibility and continuity ensure independence from the chosen $K$.
\end{proof}

Historically this branch worked on each compact $[-K,K]$ with the cone $\mathcal C_K$ generated by centered Fej\'er$\times$heat atoms $\Phi_{B,t,0}$ at the fixed scale $t=t_{\mathrm{crit}}$. The uniform Archimedean floor $c_*>0$ comes from Lemma~\ref{lem:uniform-arch-floor}, the uniform prime cap $\rho(t)\le c_*/4$ from Corollary~\ref{cor:uniform-prime-cap}, and the Lipschitz constant $L_Q(K)$ from A2. Section~\ref{sec:T5-compact} records the resulting grid-lift inequality and the centered-cone positivity statement. As explained in Remark~\ref{rem:a1-centered-obstruction}, this centered cone is not dense in $W_K$, so the old last step to the full Weil class is no longer part of the present proof.
